\section{A Hoare logic}
\label{sec:hoare}

\begin{definition}[assertions]
\label{def:assn}
An \emph{assertion} is a predicate on stores, that is, a map
$P : \Stores \to \mathrm{Prop}$. We write $\Assn$ for the assertions and use
$P, Q, R$ for them throughout. Assertions are semantic objects; no assertion
language is fixed, because none is needed for soundness.
\end{definition}

\begin{definition}[partial correctness triple]
\label{def:triple}
For $P, Q \in \Assn$ and a command $c$, the triple $\triple{P}{c}{Q}$ holds
when
\[
  \forall \sigma, \tau .\;
  P(\sigma) \;\wedge\; \bigstep{c}{\sigma}{\tau} \;\Rightarrow\; Q(\tau) .
\]
This is \emph{partial} correctness: a $c$ that never terminates satisfies every
triple, since \Cref{def:bigstep} produces no derivation for it.
\end{definition}

\begin{definition}[assertion substitution]
\label{def:subst}
For $P \in \Assn$, a variable $x$ and an expression $a$, the assertion
$P[a/x]$ is
\[
  P[a/x](\sigma) \;=\; P(\upd{\sigma}{x}{\sem{a}\sigma}) .
\]
Because assertions are semantic (\Cref{def:assn}), substitution is composition
with an update rather than a syntactic operation.
\end{definition}

\begin{lemma}[skip rule]
\label{lem:hoare-skip}
$\triple{P}{\skipk}{P}$ holds for every $P$.
\end{lemma}

\begin{proof}
The only rule of \Cref{def:bigstep} concluding $\bigstep{\skipk}{\sigma}{\tau}$
forces $\tau = \sigma$.
\end{proof}

\begin{lemma}[assignment rule]
\label{lem:hoare-assign}
$\triple{P[a/x]}{x := a}{P}$ holds for every $P$, $x$ and $a$.
\end{lemma}

\begin{proof}
Inverting the assignment rule of \Cref{def:bigstep} gives
$\tau = \upd{\sigma}{x}{\sem{a}\sigma}$, which is exactly the store at which
\Cref{def:subst} evaluates $P$.
\end{proof}

\begin{lemma}[sequencing rule]
\label{lem:hoare-seq}
If $\triple{P}{c_1}{Q}$ and $\triple{Q}{c_2}{R}$ then $\triple{P}{c_1;c_2}{R}$.
\end{lemma}

\begin{proof}
Invert the sequencing rule of \Cref{def:bigstep} to obtain the intermediate
store, and apply the two hypotheses in turn.
\end{proof}

\begin{lemma}[conditional rule]
\label{lem:hoare-if}
If $\triple{P \wedge b}{c_1}{Q}$ and $\triple{P \wedge \lnot b}{c_2}{Q}$, then
\[
  \triple{P}{\ifk\ b\ \thenk\ c_1\ \elsek\ c_2}{Q} ,
\]
where $P \wedge b$ denotes $\sigma \mapsto P(\sigma) \wedge \sem{b}\sigma =
\mathsf{tt}$.
\end{lemma}

\begin{proof}
Two cases, one per conditional rule of \Cref{def:bigstep}; each carries exactly
the hypothesis on $\sem{b}\sigma$ that the corresponding premise needs.
\end{proof}

\begin{lemma}[loop rule]
\label{lem:hoare-while}
If $\triple{P \wedge b}{c}{P}$ then
$\triple{P}{\whilek\ b\ \dok\ c}{P \wedge \lnot b}$.
\end{lemma}

\begin{proof}
Induction on the derivation of $\bigstep{\whilek\ b\ \dok\ c}{\sigma}{\tau}$,
generalising over $\sigma$. The $\mathsf{ff}$ case is immediate; the
$\mathsf{tt}$ case applies the hypothesis to the body and the induction
hypothesis to the remaining iterations.
\end{proof}

\begin{proposition}[rule of consequence]
\label{prop:hoare-conseq}
If $P' \Rightarrow P$ pointwise, $\triple{P}{c}{Q}$, and $Q \Rightarrow Q'$
pointwise, then $\triple{P'}{c}{Q'}$.
\end{proposition}

\begin{proof}
Immediate from \Cref{def:triple}: weaken the hypothesis before applying the
triple and strengthen its conclusion afterwards.
\end{proof}
